% Замените поясняющий текст собственными результатами.
% В шаблоне нет готового решения и обязательных внешних изображений.
\reportelement{Введение}

Подключаем к приложению PostgreSQL, учимся хранить данные вне контейнера и разбираемся с сетями Docker.~\ref{tab:parameters}

\begin{table}[H]
  \caption{Персональные параметры}\label{tab:parameters}
  \centering
  \begin{tabularx}{\textwidth}{|p{0.46\textwidth}|X|}
    \hline
    Параметр & Значение\\ \hline
    Номер ИСУ & \StudentISU\\ \hline
    Персональный порт (Postgres) & 5462\\ \hline
    Персональный порт (Приложение ассистент) & 8032\\ \hline
    Секретный ключ & secret32karpov\\ \hline
    Имя базы данных & karpov32\\ \hline
    Персональный префикс & karpov\\ \hline
    Окружение & linux/x86\_64\\ \hline
  \end{tabularx}
\end{table}

Результат хранится в ветке \texttt{lab-\LabNumber} форка:
\url{\RepositoryURL}. 

\chapter{Запуск контейнера с именованным volume}

Собираем приложения и образ postgres (рис. 1):
\reportscreenshot{1.png}{Результат проверки работы curl}{fig:proof1}

На рисунке~\ref{fig:proof1} видно что в ответе заметки добавились.


\chapter{Тест персистентности: пересоздание контейнера}
Запишем данные, пересоздадим контейнер и проверим, что данные сохранились в базу(рис. ~\ref{fig:proof2}): 

\reportscreenshot{2.png}{Артефакт персистентности}{fig:proof2}

Как видно ,после удаления контейнера и создания заново данные существуют.


\chapter{Bind mount для конфигурации}

Запустим контейнер с биндингом и проверяем содержимое конфига внутри контейнера:
\reportscreenshot{3.png}{Содержимое конфига}{fig:proof4}
Видно что контейнер выводит нужное содержимое конфига. Попробуем изменить файл на хосте, перезапустить контейнер и проверить снова содержимое изнутри контейнера:
\reportscreenshot{4.png}{Перезапуск контейнера с измененнным файлом}{fig:proof5}

Как видно, данные в конфиге поменялись

\chapter{Пользовательские сети для изоляции}

Создадим сеть karpov-backend и и пересоздадим контейнеры, подключив их к сети. Потом пингуем с 2х контейнеров в сети и вне сети:
\reportscreenshot{5.png}{Создаем сеть}{fig:proof6}

\chapter{Финальные артефакты}
Финальные артефакты работы:

\reportscreenshot{6.png}{Создаем сеть}{fig:proof7}






\reportelement{Заключение}
В результате выполнения работы удалось собрать образы, настроить базу данных и сети.

% Оставьте только использованные источники и цитируйте их в тексте.
\begin{thebibliography}{9}
  \bibitem{docker-docs}
  Docker Documentation: документация Docker.
  URL: \url{https://docs.docker.com/}
  (дата обращения: 09.11.2026).
\end{thebibliography}

